\section{A necessary pair-genus and partner-support gate for tail two}
\label{sec:factorwise-two-tail-gate}

\paragraph{Status and scope.}
HD-FT13--15, reserved question Q2dm, is a new derivation by
\texttt{/root/g34\_eightfolds}.  The elementary estimates are effective
for \(n\ge64\).  The final branch-list consequence uses the published
but unevaluated constant \(N_{14}\) of Monderer--Neftin; no numerical
value is assigned to it.  This is a necessary gate, not a complete
exclusion or a claim of an exhaustive catalogue of tensor partners.
Complete author and separate root whole-proof readings pass. Root also read the complete checker, first output and source/status record, the full invoked published proof algorithm, and visually inspected the entire branch table. No new root execution or numerical N14 is claimed.

Use the actual quotient and factorization of HD-FT01.  In particular,
there is a connected finite cover of the sphere with an actual quotient
onto \(P=A_n\) or \(S_n\), and an actual coefficient factorization
\(V=W\otimes U\), up to scalar operators.  Write \(x_i\in P\) for
the projected permutations of the genuine inertia operators \(T_i\).
The hypotheses used below are precisely
\[
 a=\dim W,\qquad b=\dim U,\qquad e=ab,\qquad r\ge12,
 \qquad \sum_{i=1}^r\operatorname{rank}(1-T_i)=2e+6.
 \tag{FT2P.1}
\]
Every \(T_i\ne1\), the tuple has product one and generates the actual
covering group.  The inequality \(r\ge12\) belongs to the original
rank-nine inducing sphere source.  It is not assigned to the projected
\(A_n\)- or \(S_n\)-factor.  The degree, coefficient field and all
scalar phases of \(U\) are arbitrary.

Let \(W\) have label \((n-2,2)\) or \((n-2,1,1)\), its transpose,
or the irreducible ordinary restriction of one of these representations
to \(A_n\).  The restrictions are irreducible for \(n\ge64\).
All comparisons below are absolute-character or projective-support
comparisons, so sign twists have exactly the same bounds.

\subsection*{HD-FT13. An all-eigenvalue comparison with unordered pairs}

For a permutation \(x\), put \(f=X_1(x)\), \(u=X_2(x)\), \(m=n-f\)
and
\[
 B=\binom n2,\qquad F=\binom f2+u.
 \tag{FT2P.2}
\]
Thus \(F\) is the number of fixed unordered pairs.  The exact two
characters, in the order of the two labels above, are
\[
 \begin{array}{c|c|c}
 &\chi(x)&a\\ \hline
 R&F-f=f(f-3)/2+u&B-n,\\
 H&F-f+1-2u=(f-1)(f-2)/2-u&B-n+1.
 \end{array}
 \tag{FT2P.3}
\]
These are the tail-two rows of the already reviewed character-polynomial
table.\cite[Table~5.1, p.~31; Proposition~I.1 and proof]{HD:GarsiaGoupil2009}
They also follow from the unordered-pair permutation representation and
the exterior square of the deleted permutation representation.

\begin{lemma}[HD-FT13]
For every \(n\ge64\) and every permutation,
\[
 a\ge n^2/4,\qquad
 \frac{|\chi(x)|}{a}\le \frac FB+
                      \frac8{n^2}\,\mathbf1_{f\le2}.
 \tag{FT2P.4}
\]
Let \(I_2(x)\) be the index of \(x\) on unordered pairs and
\(I(x)\) its natural index.  If
\(\nu(Wx)=a-\max_\zeta\dim\ker(Wx-\zeta)\), then
\[
 \frac{\nu(Wx)}a\ge\frac{I_2(x)}B-
                  \frac8{n^2}\,\mathbf1_{f\le2},\qquad
 \frac{\nu(Wx)}a\ge\frac{I(x)}n-
                  \frac1{n-1}-\frac8{n^2}\,\mathbf1_{f\le2}.
 \tag{FT2P.5}
\]
Both conclusions concern every target eigenvalue, including arbitrary
scalar twists of \(Wx\).
\end{lemma}

\begin{proof}
The dimension bounds and \(B\ge n^2/4\) follow immediately from
(FT2P.3).  Also \(2u\le m\).  For the first character,
\[
 B\chi_R-(B-n)F
       =\frac n2(-fm+2u).
 \tag{FT2P.6}
\]
If \(f\ge1\), the last parenthesis is at most \(m(1-f)\le0\).
If \(f=0\), the normalized difference is
\(nu/((B-n)B)\le8/n^2\).
The only possible negative values of \(\chi_R\) occur at \(f=1,2\),
and are at least \(-1\).  Their absolute normalized value is at most
\(1/a\le4/n^2\).  This proves (FT2P.4) for \(R\), with no error
when \(f\ge3\).

For the hook character the exact signed comparison is
\[
 B\chi_H-(B-n+1)F
   =-\frac{n-1}{2}\bigl(m(f-1)+2(n-1)u\bigr).
 \tag{FT2P.7}
\]
It is nonpositive for \(f\ge1\).  For \(f=0\) its numerator is
at most \(B\), so its normalized positive error is at most
\(1/a\le4/n^2\).  Negative character values are handled separately,
using \(-\chi_H=F-(f-1)^2\):
\[
 B(-\chi_H)-(B-n+1)F
   =(n-1)F-B(f-1)^2
   \le\frac{(n-1)^2}{2}f(2-f).
 \tag{FT2P.8}
\]
Here the inequality uses \(2u\le n-f\).  The right side is
nonpositive at \(f=0\) and every \(f\ge2\).  At \(f=1\), division
by \(aB\ge n^4/16\) bounds the error by \(8/n^2\).
This proves the asserted absolute comparison in all cases.

Every power of \(x\) fixes all the letters fixed by \(x\).  Thus
if \(f(x)\ge3\), all powers have the error-free estimate.
For \(d=\operatorname{ord}(x)\), cyclic Fourier inversion gives,
for every \(\zeta\),
\[
 \frac{\dim\ker(Wx-\zeta)}a
 \le\frac1d\sum_{j=0}^{d-1}\frac{|\chi(x^j)|}a
 \le\frac1d\sum_{j=0}^{d-1}\frac{F(x^j)}B
                       +\frac8{n^2}\mathbf1_{f(x)\le2}.
 \tag{FT2P.9}
\]
The average of \(F(x^j)\) is the number of cycles on unordered
pairs, proving the first support estimate.
Finally,
\[
 \frac FB\le\frac fn+\frac1{n-1}.
 \tag{FT2P.10}
\]
Indeed for \(f\ge1\), the stronger inequality \(F/B\le f/n\)
follows from
\(2F-(n-1)f\le(f-1)(f-n)\le0\); for \(f=0\) use
\(F/B\le1/(n-1)\).  Averaging (FT2P.10) over powers and using
\(d^{-1}\sum_j f(x^j)=n-I(x)\) proves the second estimate.
\end{proof}

Spectral convolution with any finite-order partner, and with any
scalar phase, gives
\[
 \frac{\operatorname{rank}(1-T_i)}e\ge\frac{\nu(Wx_i)}a.
 \tag{FT2P.11}
\]
For example, each partner eigenspace can contribute at most the
largest \(W\)-eigenspace to a specified product eigenvalue.  This
explains why no reality or coefficient-field equality is needed.

\subsection*{HD-FT14. At most four error-bearing entries and pair genus at most fourteen}

\begin{proposition}[HD-FT14]
In the actual source (FT2P.1), at most four projected inertia
permutations satisfy \(f(x_i)\le2\).  The connected quotient cover
on unordered pairs has genus \(g_2\le14\).
\end{proposition}

\begin{proof}
Call an entry with \(f(x_i)\le2\) bad.  Its natural index satisfies
\(I(x_i)\ge(n-2)/2\), since every moved cycle has length at least two.
Equations (FT2P.5) and (FT2P.11) imply
\[
 \frac{\operatorname{rank}(1-T_i)}e
 \ge\frac12-\frac1n-\frac1{n-1}-\frac8{n^2}
 >\frac9{20}\qquad(n\ge64).
 \tag{FT2P.12}
\]
For the strict inequality it suffices to check
\(1/64+1/63+1/512<1/20\), since each subtracted term decreases.
On the other hand,
\[
 \frac{2e+6}{e}\le2+\frac{24}{n^2}<\frac94.
 \tag{FT2P.13}
\]
Five bad entries would contradict these two strict inequalities.

The unordered-pair action of \(P\) is transitive.  Hence its actual
degree-\(B\) quotient is connected and Riemann--Hurwitz gives
\(\sum_i I_2(x_i)=2B-2+2g_2\), including zero contributions from
projected identity entries.  Summing the first inequality of
(FT2P.5), with at most four errors, yields
\[
 2+\frac6e\ge2+\frac{2g_2-2}{B}-\frac{32}{n^2},
 \qquad
 2g_2-2\le\frac{56B}{n^2}
             =28\left(1-\frac1n\right)<28.
 \tag{FT2P.14}
\]
The integral genus is therefore at most fourteen.
\end{proof}

\subsection*{HD-FT15. The actual branch list forces a small-support partner operator}

Let \(N_{14}\) denote the constant in Monderer--Neftin Theorem~1.2
for genus bound fourteen.  This is an existing published qualitative
constant; its value is not evaluated here.
\cite[Theorems~1.2 and 5.1; Proposition~4.1 and Table~2.1]{HD:MondererNeftin2022}

\begin{theorem}[HD-FT15]
Suppose \(n\ge64\) and \(n>N_{14}\).  In every source (FT2P.1),
the natural and unordered-pair quotient covers both have genus zero.
At most five \(x_i\) are nonidentity, so at least seven genuine
original inertia operators have \(x_i=1\).  They satisfy
\[
 \sum_{x_i=1}\frac{\operatorname{rank}(1-T_i)}e
       \le\frac{61}{n^2}.
 \tag{FT2P.15}
\]
In particular at least one such genuine operator has
\[
 0<\frac{\operatorname{rank}(1-T_i)}e\le\frac{61}{7n^2},
 \qquad b\ge\frac{7n^2}{61}.
 \tag{FT2P.16}
\]
If these kernel-projected genuine operators all have normalized moving
rank at least a specified \(\eta>0\), then necessarily
\[
                      7\eta\le\frac{61}{n^2}.
 \tag{FT2P.17}
\]
This last assertion is conditional on the stated actual support floor;
no such uniform floor for an unbounded partner family is assumed.
\end{theorem}

\begin{proof}
Apply Monderer--Neftin Theorem~1.2 to the actual natural degree-\(n\)
cover and its pair-stabilizer subgroup \(H=P\cap(S_{n-2}\times S_2)\).
The subgroup \(H\) cannot contain \(A_{n-1}\), already by its order
when \(n\ge64\).  Proposition HD-FT14 supplies the genus bound
fourteen required by that theorem.  It follows that the natural
ramification is among the types in Proposition~4.1 and Table~2.1.
Proposition~4.1(1), for either \(A_n\) or \(S_n\), gives genus zero
for the point and pair stabilizers.  Every row of Table~2.1 has at
most five genuine branch points: only rows F1.1 and F1.2 have five,
and all others have at most four.  Multiplicities such as ``twice''
and ``thrice'' count distinct branch points.  These are the projected
nonidentity entries, not the entire original tuple.  Since \(r\ge12\),
at least seven original genuine entries project to the identity.

The support comparison and nonnegative pair genus give
\[
 \sum_{x_i\ne1}\frac{\nu(Wx_i)}a
 \ge2-\frac2B-\frac{32}{n^2}
 \ge2-\frac{37}{n^2}.
 \tag{FT2P.18}
\]
The last inequality uses \(2/B=4/(n(n-1))\le5/n^2\), valid for
\(n\ge5\).  Subtract (FT2P.18), via (FT2P.11), from the actual
budget (FT2P.1), and use \(6/e\le24/n^2\).  This proves
(FT2P.15) and its seven-entry average in (FT2P.16).

If \(x_i=1\), the actual operator is \(1_W\otimes V_i\) after
absorbing its scalar in \(V_i\).  Since \(T_i\ne1\), one has
\(V_i\ne1\) and
\(\operatorname{rank}(1-T_i)/e
=\operatorname{rank}(1-V_i)/b\ge1/b\).
Summing over the seven entries proves the partner-dimension bound.
The same sum proves (FT2P.17) under its explicit hypothesis.
\end{proof}

\paragraph{Retained boundary.}
The result does not exclude an unbounded tensor partner that supplies
such small-support genuine operators.  In particular it does not
identify a maximal-degree ordinary factor, classify factor permutations,
or establish that all alternative partner images are in a bounded
catalogue.  Any consequence requiring those statements must prove them
separately.  The effective estimates through HD-FT14 and the qualitative
\(N_{14}\)-dependent conclusion of HD-FT15 are deliberately distinct;
the numerical-source question Q2dk remains open.
